\section{Phase 2: Classifier Benchmarking}
\label{sec:phase2}
\label{sec:phase2_method}

\subsection{Objectives}
\label{sec:phase2_objectives}

Following the separability-guided optimization of Phase~1, the second phase
evaluates the predictive performance of multiple machine learning algorithms on
the optimized feature representations. The goal is not merely to identify the
single best classifier, but to understand how different learning paradigms
interact with the various handcrafted feature sets. To this end, eight
representative algorithms---spanning linear, distance-based, tree-based,
ensemble, boosting, and neural-network approaches---are benchmarked under
identical experimental conditions, providing a broader and fairer comparison
than the one- or two-classifier evaluations common in prior work.

\subsection{Benchmarking Methodology}
\label{sec:phase2_methodology}

The benchmarking framework was designed for reproducibility, fairness, and
statistical rigor. For each feature representation of
Section~\ref{sec:features}, the following steps were applied: feature cleaning
and standardization; hyperparameter optimization by grid search; model fitting
with stratified cross-validation; recording of performance metrics; and
statistical validation. This procedure was repeated for every
classifier--feature-set combination, giving the pipeline optimized features
$\rightarrow$ cleaning $\rightarrow$ standardization $\rightarrow$ hyperparameter
search $\rightarrow$ cross-validation $\rightarrow$ testing $\rightarrow$
statistical validation.

\subsection{Evaluated Classifiers}
\label{sec:phase2_classifiers}

Eight algorithms were selected to ensure diversity in learning strategy.
\emph{Support Vector Machine} (SVM)~\cite{cortes1995} seeks a maximum-margin
separating hyperplane and is effective in high-dimensional spaces, with
nonlinear modeling through kernel functions. \emph{K-Nearest Neighbors} (KNN)
is a non-parametric method that assigns each sample the majority label among its
$k$ nearest neighbors, providing a useful baseline for feature-space
separability. \emph{Random Forest} (RF)~\cite{breiman2001} aggregates decision
trees trained on bootstrap samples and random feature subsets, offering
robustness, resistance to overfitting, and built-in feature-importance
estimation. \emph{XGBoost}~\cite{chen2016} builds gradient-boosted trees in which
each tree corrects its predecessors' errors, combining high accuracy with
regularization. \emph{LightGBM}~\cite{ke2017} is an efficient gradient-boosting
framework using histogram-based, leaf-wise tree growth for fast, memory-efficient
training. \emph{Logistic Regression} (LR) is an interpretable linear
probabilistic baseline. The \emph{Multi-Layer Perceptron} (MLP) is a feed-forward
network that operates directly on handcrafted feature vectors, bridging classical
and neural approaches. Finally, \emph{Extra Trees} (ET) resembles RF but
introduces additional randomness in node splitting, increasing tree diversity and
reducing variance.

\subsection{Hyperparameter Optimization}
\label{sec:phase2_hyperparams}

Classifier performance depends strongly on parameter selection, so hyperparameter
optimization was performed independently for each classifier using grid search
with stratified cross-validation, selecting the configuration with the highest
macro-$F_1$. Where supported, the class weighting was also tuned
(\texttt{balanced} vs.\ uniform) to counter the class imbalance.
Table~\ref{tab:hparams} summarizes the search spaces.

\begin{table}[t]
\centering
\caption{Hyperparameter search spaces explored by grid search.}
\label{tab:hparams}
\renewcommand{\arraystretch}{1.25}
\footnotesize
\begin{tabular}{@{}ll@{}}
\toprule
\textbf{Classifier} & \textbf{Search space} \\
\midrule
SVM      & kernel $\in\{$linear, rbf, poly, sigmoid$\}$; \\
         & $C \in \{10^{-3},\dots,10^{3}\}$; $\gamma \in \{$scale, auto, $10^{-4}\text{--}10\}$ \\
KNN      & $k \in \{1,2,3,5,\dots,51\}$; \\
         & metric $\in\{$euclidean, manhattan, chebyshev, minkowski$\}$; \\
         & weights $\in\{$uniform, distance$\}$ \\
RF       & trees $\in\{50,\dots,1000\}$; depth $\in\{$None$,5,\dots,80\}$; \\
         & min.\ split $\in\{2,5,10,20\}$ \\
ET       & trees $\in\{100,300,500\}$; min.\ split $\in\{2,5,10\}$ \\
XGBoost  & trees $\in\{100,\dots,800\}$; depth $\in\{3,\dots,8\}$; \\
         & learning rate $\in\{0.01,\dots,0.2\}$; $\gamma \in\{0,0.1,0.3\}$ \\
LightGBM & trees $\in\{100,\dots,800\}$; depth $\in\{-1,4,6,8\}$; \\
         & learning rate $\in\{0.01,\dots,0.2\}$; leaves $\in\{15,31,63,127\}$ \\
LR       & $C$ (inverse regularization strength) \\
MLP      & layers $\in\{(256),(512),(256,128),(512,256),(512,256,128)\}$; \\
         & activation $\in\{$relu, tanh$\}$; $\alpha \in \{10^{-4},10^{-3},10^{-2}\}$ \\
\bottomrule
\end{tabular}
\end{table}

\subsection{Cross-Validation Strategy}
\label{sec:phase2_cv}

To obtain robust, unbiased performance estimates, stratified five-fold
cross-validation was used within the training partition, preserving class
proportions in each fold. In each iteration, four folds are used for training and
one for validation, and metrics are averaged across folds. This reduces the
influence of any single data partition relative to a lone train--test split.

\subsection{Evaluation Metrics}
\label{sec:phase2_metrics}

Performance was measured with accuracy and macro-averaged precision, recall, and
$F_1$ score. For a class treated as positive,
\begin{equation}
\mathrm{Accuracy} = \frac{TP + TN}{TP + TN + FP + FN},
\end{equation}
\begin{equation}
\mathrm{Precision} = \frac{TP}{TP + FP}, \qquad
\mathrm{Recall} = \frac{TP}{TP + FN},
\end{equation}
\begin{equation}
F_1 = 2 \cdot \frac{\mathrm{Precision}\cdot \mathrm{Recall}}
{\mathrm{Precision} + \mathrm{Recall}},
\end{equation}
where $TP$, $TN$, $FP$, and $FN$ denote true/false positives and negatives.
Because the task has three classes, precision, recall, and $F_1$ are
macro-averaged so that each tumor class contributes equally regardless of its
prevalence.

\subsection{Benchmark Design}
\label{sec:phase2_design}

The study evaluates the four individual feature families (A, B, C, D) and the
four fusion sets (A+B, B+C, A+B+C, Full(opt)) across all eight classifiers,
producing a full classifier--feature-set comparison matrix. This design directly
supports the research questions on the most discriminative feature representation
and the most effective classifier (RQ2 and RQ3), which are stated in full and
answered in Section~\ref{sec:results}.

\subsection{Summary}
\label{sec:phase2_summary}

Phase~2 establishes a reproducible benchmarking framework combining systematic
hyperparameter optimization, stratified cross-validation, multiple evaluation
metrics, and broad classifier diversity. The resulting experimental findings and
comparative analyses are presented in the following section.
